\clearpage
\chapterepigraph{All things are implicated with one another, and the bond is
holy; and there is hardly anything unconnected with any other thing.}{Marcus
Aurelius}{\textit{Meditations}}

\begingroup
\let\clearpage\relax
\let\cleardoublepage\relax
\chapter{Conclusions}
\label{ch:conclusions}
\endgroup

The work formed a feedback cycle between solver and physics. We constructed a variational
state and compared it with external references wherever they were available; the energies
agree closely. We then used the state as an instrument to resolve the Wigner regime: rings at
definite occupancies, $n$-fold angular order that becomes large only in the last decade of
confinement, and a shell radius approaching its classical $\omega^{-2/3}$ scaling. Those are the
measurements a sampler needs in order to place its mass correctly. Feeding the measured
geometry back into the proposal carried a cascade with MCMC-free updates from
$\omega\approx0.01$, where the fixed origin-centred sampler fails, to a stable checkpoint at
$\omega=0.0035$.

\emph{Architecture.} The implemented model families differ before their energies do
(\S\ref{sec:kernel-q1}). Under the parametrisations used here, the message-passing correlator
concentrates its Fisher weight in fewer effective directions and is better described by the
selected operator dictionary. Its lower local-energy variance is a state-level difference.
The correlators also use different primitive features, so the present comparison does not
attribute that difference to routing alone. At weak confinement the conventional backflow
falls to rank one and loses half a percent of the energy, but it is not matched to the CTNN in
capacity, input scaling or centre-of-mass handling. This is an implementation result until
those controls are supplied. The correlator swap moves the same boundary and is consistent
with, rather than proof of, a shared relational bottleneck.

\emph{Optimiser.} Conditioning alone does not decide whether natural-gradient
preconditioning earns its cost (\S\ref{sec:kernel-cond}). Under $|\Psi|^2$ sampling the
empirical metric is badly conditioned on the very models where Adam and stochastic
reconfiguration agree to four decimals. Under a fixed collocation proposal, SR is required
in the tested configurations and becomes unstable after proposal overlap collapses. We read
that pattern as saying that estimability, rather than conditioning,
separates the regimes. It is a hypothesis the failures support, not yet a criterion: the
metric was never measured under the collocation measure itself.

\emph{Paradigm.} An accurate wavefunction can be trained with no Markov chain in the
gradient, inside a regime whose edges we found by crossing them
(\S\ref{sec:kernel-mcmcfree}). The claim is about the gradient only: a Metropolis probe still
selects checkpoints and certifies the final energy, and the cascades start from
reference-guided pretrained states. The weak-form objective comes within about a tenth of a
percent of DMC at $N\in\{6,12\}$ under strong confinement and within about a third of a
percent of our own variational energies well into the correlated regime, as measured by
independent Metropolis evaluation of every trained state. The failures are consistent with
proposal mismatch, heavier weights, gradient noise and an unreliable empirical metric, but
the interventions do not isolate every link. A better proposal moves that limit. The cascade then meets a second
failure, of transfer between confinements; trained directly from a pretrained checkpoint
instead, the same ansatz form reaches $\omega=10^{-3}$.

\emph{Physics.} The trained state serves as an instrument
(\S\ref{sec:wigner-molecule}). The crossover it resolves is staged: the cloud expands and its
shells stiffen radially well before the electrons acquire angular order within them. The
shell geometry that emerges at the weakest traps is resolved topology by topology, and its
measured radius approaches the classical value $1.334\,\omega^{-2/3}$ to within $1.7\%$. The
staging itself is better established than where the transition sits: a grid with one point
per decade can order the stages but not locate them. And these are the lowest states found
at fixed spin projection, in the regime where a change of ground-state spin is most
plausible.

\section{Outlook}
\label{sec:outlook}

Three papers based on this work are in preparation. The steps that would strengthen the
method itself are set out in \S\ref{sec:implications}; what follows concerns the physics the
same machinery could reach. Each direction changes the Hamiltonian or the question while
keeping the variational principle, the optimiser and diagnostics that act only on samples of
$|\Psi|^2$.

The most immediate is spin. Every state here is the lowest found at fixed spin projection,
and a scan over total spin would settle whether the ground state changes spin in the deep
Wigner regime, where the structural conclusions are strongest. Larger $N$ follows once the
quadratic memory cost of the backflow is removed, and would test the shell sequences of
classical clusters across several filled shells.

The confinement itself can be changed. A multi-well potential, as in coupled
dots~\cite{Mariadason2018-DoubleDotHF} and the spin qubits built from them, enters the
Hamiltonian directly, but it breaks the rotational symmetry on which both the shell diagnostics and the
collocation proposal rely, and both would have to be rebuilt around several centres. A
perpendicular magnetic field breaks time-reversal symmetry instead. The harmonic-oscillator
orbitals become Fock--Darwin orbitals, the ansatz needs a phase as well as an amplitude, and
the ground state passes through a sequence of angular-momentum transitions as the field
grows. The Wigner crossover under a field is a natural test
for the structural diagnostics developed here.

Two further directions change the question rather than the system. The quantum geometric
tensor studied in Chapter~\ref{ch:kernel} is the metric that the time-dependent variational
principle inverts, so real-time evolution of a neural state~\cite{Carleo2017-NQS}, such as
the response of a Wigner molecule to a sudden change of the trap, would inherit the
conditioning and estimability questions raised here directly. And Rényi entanglement
entropies between regions or shells of the dot can be estimated from samples of two
independent copies of the trained state, giving the crossover a measure that neither the
energy nor the density provides.

\section{Concluding perspective}

All four answers are bounded in the same way. Every state is the lowest found within a fixed
spin-projection sector; the architecture comparison is seeded at $N{=}6$ and thinner above
it, and its tangent-space measures belong to the networks as parametrised. The backflow
variants are not matched in capacity, input scaling or centre-of-mass handling, and the
correlators use different primitive features, so neither comparison isolates message passing.
\S\ref{sec:implications} sets out what would move
each of those.

For the well-resolved $N\le12$ studies, the observed failures were more closely associated
with alignment, conditioning and sampling than with raw parameter count. This does not make
capacity irrelevant. At $N{=}20$ the memory cost of the backflow's edge tensor forces a
capacity cut and coincides with degraded collocation gradients. Across architecture, optimisation and
sampling, performance is repeatedly associated with alignment to the collective structure of
the problem: breathing, correlation-hole and relative modes at strong confinement, and the
shell geometry of the Wigner molecule at weak confinement. Whether this is one mechanism or
several related ones remains open.

Several of these results were reached only after an initial reading of the same data proved
wrong: a rank collapse that survived seeding where a dimension inversion did not, a transfer
failure first attributed to sampler placement and later found to persist in the retained
wavefunction, an importance-sampling
error whose signature was indistinguishable from an architectural deficiency, and a
strong-confinement collocation energy that looked statistically exact until an independent
evaluation showed it to be the lucky tail of a best-of-three selection. Each was
identified by a diagnostic that could separate the competing explanations. That is the
methodological point these chapters have tried to make concrete: a variational calculation is
only as trustworthy as the tests that can tell its failure modes apart.
